\documentclass[../../../include/open-logic-section]{subfiles}

\begin{document}

\olfileid[pt]{sth}{choice}{vitali}
\olsection{Apêndice: Paradoxo de Vitali}

Para ter uma ideia real de se a construção de Banach--Tarski é
aceitável ou não, devíamos examinar a sua \emph{prova}. Infelizmente,
isso exigiria muito mais álgebra do que podemos apresentar aqui.
Contudo, podemos oferecer algumas observações rápidas que talvez esclareçam
a prova de Banach--Tarski,\footnote{Para um tratamento bem mais completo,
ver \cite{Weston2003} ou \cite{Wagon2016}.} concentrando-nos no
resultado seguinte:

\begin{thm}[Paradoxo de Vitali (em $\ZFC$)]\ollabel{vitaliparadox}
Qualquer círculo pode ser decomposto numa quantidade enumerável de peças, que podem ser
remontadas (por rotação e transporte) de modo a formar duas cópias desse
círculo.
\end{thm}

O Paradoxo de Vitali é muito mais fácil de provar do que o Paradoxo de Banach--Tarski. Chamamos-lhe ``Paradoxo de Vitali'', pois segue-se da construção, em
\citeyear{Vitali1905}, de um conjunto não mensurável. Mas os
aspectos conjuntistas da prova do Paradoxo de Vitali e do
Paradoxo de Banach--Tarski são muito semelhantes. A diferença essencial
entre os resultados é que Banach--Tarski considera uma
decomposição \emph{finita}, enquanto o Paradoxo de Vitali considera uma
decomposição \emph{enumerável infinita}. Como \citet{Weston2003}
diz, o Paradoxo de Vitali ``certamente não é tão notável como o
paradoxo de Banach--Tarski, mas ilustra que paradoxos geométricos podem
ocorrer mesmo em situações `simples'.'' 

O Paradoxo de Vitali trata de uma figura bidimensional, um círculo. Logo
trabalhamos no plano, $\Real^2$. Seja $\rotationsgroup$ o conjunto das
rotações (no sentido dos ponteiros do relógio) de pontos em torno da origem por valores em radianos \emph{racionais}
no intervalo $[0,2\pi)$. Eis alguns fatos algébricos sobre
$\rotationsgroup$ (se não compreender o enunciado do
resultado, a prova esclarecerá o seu significado):

\begin{lem}\ollabel{rotationsgroupabelian}
$\rotationsgroup$ forma um {grupo} abeliano sob a composição de funções.
\end{lem}

\begin{proof}
Escrevendo $0_{\rotationsgroup}$ para a rotação por $0$ radianos, este é
um elemento neutro para $\rotationsgroup$, pois
$\comp{0_{\rotationsgroup}}{\rho} = \comp{\rho}{0_{\rotationsgroup}} =
\rho$ para qualquer $\rho \in \rotationsgroup$.

Todo o elemento tem inverso. Onde $\rho \in \rotationsgroup$ roda
por $r$ radianos, $\rho^{-1} \in \rotationsgroup$ roda por $2\pi - r$
radianos, de modo que $\rho \circ \rho^{-1} = 0_\rotationsgroup$.

A composição é associativa: $\comp{\rho}{(\comp{\sigma}{\tau})} =
\comp{(\comp{\rho}{\sigma})}{\tau}$ para quaisquer $\rho, \sigma, \tau \in
\rotationsgroup$

A composição é comutativa: $\comp{\rho}{\sigma} =
\comp{\sigma}{\rho}$ para quaisquer $\rho, \sigma \in \rotationsgroup$.
\end{proof}

De fato, podemos partir o nosso grupo $\rotationsgroup$
ao meio e depois usar qualquer metade para recuperar o grupo inteiro:

\begin{lem}\ollabel{disjointgroup}
Existe uma partição de $\rotationsgroup$ em dois conjuntos disjuntos,
$\rotationsgroup_{1}$ e $\rotationsgroup_{2}$, ambos dos quais são uma
base para $\rotationsgroup$. 
\end{lem}

\begin{proof}
Seja $\rotationsgroup_{1}$ o conjunto das rotações por valores racionais em radianos
em $[0, \pi)$; seja $\rotationsgroup_{2} = \rotationsgroup
\setminus \rotationsgroup_1$. Por álgebra elementar,
$\Setabs{\comp{\rho}{\rho}}{\rho \in \rotationsgroup_1} =
\rotationsgroup$. Um resultado semelhante obtém-se para
$\rotationsgroup_2$.
\end{proof}

Usaremos este fato sobre grupos para estabelecer \olref{vitaliparadox}.
Seja $\onesphere$ o círculo unitário, isto é, o conjunto dos pontos
exatamente a~$1$ unidade da origem do plano, isto é,
$\Setabs{\tuple{r,s} \in \Real^2}{\sqrt{r^2+s^2}=1}$. Partiremos
$\onesphere$ em partes considerando a relação seguinte em
$\onesphere$:
\[
	r \sim s \emph{ sse }(\exists \rho \in \rotationsgroup)\rho(r) = s.
\]
Ou seja, os pontos de $\onesphere$ estão ligados por esta relação sse se pode passar de um para o outro por uma rotação de valor racional em torno da origem. Sem surpresa:

\begin{lem}
$\sim$ é uma relação de equivalência.
\end{lem}

\begin{proof}
Trivial, usando \olref{rotationsgroupabelian}. 
\end{proof}

Invocamos agora a Escolha para obter um conjunto, $C$, contendo exatamente um
membro de cada classe de equivalência de $\onesphere$ sob $\sim$. Ou seja,
consideramos uma função de escolha $f$ no conjunto das classes de
equivalência,\footnote{Como $\rotationsgroup$ é !!{enumerable},
cada !!{element} de $E$ é !!{enumerable}. Como $\onesphere$ é
!!{nonenumerable}, segue-se de \olref{vitalicover} e
\olref[card-arithmetic][simp]{kappaunionkappasize} que $E$ é
!!{nonenumerable}. Logo isto é um uso de Escolha \emph{não} enumerável.}
\[
	E = \Setabs{\equivrep{r}{\sim}}{r \in \onesphere},
\]
e seja $C = \ran{f}$. Para cada rotação $\rho \in \rotationsgroup$,
o conjunto $\funimage{\rho}{C}$ consiste nos pontos obtidos ao
aplicar a rotação $\rho$ a cada ponto de $C$. Os dois resultados seguintes
mostram que estes conjuntos cobrem o círculo por completo e sem
sobreposição:

\begin{lem}\ollabel{vitalicover}
$\onesphere = \bigcup_{\rho \in \rotationsgroup} \funimage{\rho}{C}$.
\end{lem}

\begin{proof}
Fixe-se $s \in \onesphere$; existe algum $r \in C$ tal que $r \in
\equivrep{s}{\sim}$, isto é, $r \sim s$, isto é, $\rho(r) = s$ para algum
$\rho \in \rotationsgroup$. 
\end{proof}

\begin{lem}\ollabel{vitalinooverlap}
Se $\rho_1 \neq \rho_2$ então $\funimage{\rho_1}{C} \cap \funimage{\rho_2}{C} = \emptyset$. 
\end{lem}

\begin{proof}
Suponhamos $s \in \funimage{\rho_{1}}{C} \cap \funimage{\rho_{2}}{C}$. Logo
$s = \rho_{1}(r_{1}) = \rho_{2}(r_{2})$ para alguns $r_{1}, r_{2} \in C$.
Daí $\rho^{-1}_2(\rho_1(r_1)) = r_{2}$, e
$\comp{\rho_1}{\rho^{-1}_2} \in \rotationsgroup$, logo $r_{1} \sim
r_{2}$. Logo $r_1 = r_2$, pois $C$ selecciona exatamente um membro de cada
classe de equivalência sob $\sim$. Logo $s = \rho_1(r_1) = \rho_2(r_1)$,
e portanto $\rho_1 = \rho_2$.
\end{proof}

Aplicamos agora os fatos algébricos anteriores ao nosso círculo:

\begin{lem}\ollabel{pseudobanachtarski}
Existe uma partição de $\onesphere$ em dois conjuntos disjuntos, $D_{1}$
e $D_{2}$, tal que $D_{1}$ pode ser particionado numa quantidade enumerável de
conjuntos que podem ser rodados de modo a formar uma cópia de $\onesphere$ (e
analogamente para $D_{2}$).
\end{lem}

\begin{proof}
Usando $\rotationsgroup_{1}$ e $\rotationsgroup_{2}$ de \olref{disjointgroup}, seja:
\begin{align*}
	D_{1} &= \bigcup_{\rho \in \rotationsgroup_1} \funimage{\rho}{C} & 
	D_{2} &= \bigcup_{\rho \in \rotationsgroup_2} \funimage{\rho}{C}
\end{align*}
Isto é uma partição de $\onesphere$, por \olref{vitalicover}, e $D_1$
e $D_2$ são disjuntos por \olref{vitalinooverlap}. Por construção,
$D_1$ pode ser particionado numa quantidade enumerável de conjuntos,
$\funimage{\rho}{C}$ para cada $\rho \in R_1$. E estes podem ser rodados
de modo a formar uma cópia de $\onesphere$, pois $\onesphere = \bigcup_{\rho \in
\rotationsgroup}\funimage{\rho}{C} = \bigcup_{\rho \in
\rotationsgroup_1}\funimage{(\comp{\rho}{\rho})}{C}$ por
\olref{disjointgroup} e \olref{vitalicover}. O mesmo raciocínio
vale para $D_2$. \end{proof}\noindent Isto implica de imediato
o Paradoxo de Vitali. Pois podemos gerar \emph{duas} cópias de
$\onesphere$ a partir de $\onesphere$, apenas partindo-o numa quantidade enumerável de
peças (os vários $\funimage{\rho}{C}$) e depois movendo-as rigidamente
(basta rodar cada peça de $D_1$, e primeiro transportar
e depois rodar cada peça de $D_2$).

Recapitulemos a estratégia de prova. Começamos com fatos algébricos sobre o grupo de rotações no plano. Usamos este grupo para
partir $\onesphere$ em classes de equivalência. Chegamos então a um
``paradoxo'', usando a Escolha para seleccionar elementos de cada classe.

Usamos exatamente a mesma estratégia para provar Banach--Tarski. A principal
diferença é que os fatos algébricos usados na prova de Banach--Tarski
são bastante mais complicados do que os usados no Paradoxo de Vitali. Mas esses fatos algébricos nada têm a ver com a
Escolha. Resumimos-os rapidamente. 

Para provar Banach--Tarski, começamos por estabelecer um análogo de
\olref{disjointgroup}: qualquer \emph{grupo livre} pode ser partido em quatro
peças que, intuitivamente, podemos ``mover'' de modo a recuperar duas cópias
do grupo inteiro.\footnote{O fato de podermos usar \emph{quatro}
peças deve-se a \cite{Robinson1947}. Para uma prova recente, ver
\citet[Theorem 5.2]{Wagon2016}. Seguimos \citet[p.~3]{Weston2003}
ao descrever isto como ``mover'' as peças do grupo.} Mostramos depois
que podemos usar duas rotações particulares em torno da origem de
$\Real^3$ para gerar um grupo livre de rotações, $F$.\footnote{Ver
\citet[Theorem 2.1]{Wagon2016}.} (Ainda sem Escolha.) Consideramos agora pontos
na superfície da esfera como ``semelhantes'' sse um pode obter-se
do outro por uma rotação em~$F$. \emph{Usamos então a Escolha} para
seleccionar exatamente um ponto de cada classe de equivalência de pontos ``semelhantes''.
Aplicando a nossa divisão de $F$ à superfície da esfera, como
em \olref{pseudobanachtarski}, partimos essa superfície em quatro peças,
que podemos ``mover'' de modo a obter duas cópias da superfície da
esfera. E isto estabelece \citep{Hausdorff1914}:

\begin{thm}[Paradoxo de Hausdorff (em $\ZFC$)] 
A superfície de qualquer esfera pode ser decomposta num número finito de peças,
que podem ser remontadas (por rotação e transporte) de modo a formar duas
cópias disjuntas dessa esfera.
\end{thm}

Ainda são precisos alguns truques algébricos para obter o Teorema completo de
Banach--Tarski (que não diz só respeito à superfície da esfera,
mas também ao seu interior). Francamente, porém, isto é apenas a cereja no topo do
bolo algébrico. Por isso Weston escreve:
\begin{quote}	
  [\ldots] o resultado sobre grupos livres é o passo \emph{chave} na
  prova do paradoxo de Banach--Tarski. Deste ponto de vista, o
  paradoxo de Banach--Tarski não é tanto uma afirmação sobre $\Real^3$ quanto
  uma afirmação sobre a complexidade do grupo [das translações
  e rotações em $\Real^3$]. \cite[p.~16]{Weston2003}
\end{quote}
Ou seja: podermos ou não oferecer uma decomposição \emph{finita} (como em
Banach--Tarski) ou uma decomposição \emph{enumerável infinita} (como no
Paradoxo de Vitali) resume-se a certos fatos de teoria dos grupos sobre
trabalhar em duas ou três dimensões.

Admitidamente, esta última observação estraga ligeiramente a piada no final
de \olref[banach]{sec}. Por ser bidimensional,
``Banach--Tarski'' tem de ser dividido numa \emph{infinidade} enumerável de
peças, se se quiserem rearranjar essas peças de modo a formar ``Banach--Tarski
Banach--Tarski''. Para consertar a piada, é preciso escrever em três
dimensões. Deixamos isto como exercício para o leitor.

Um último comentário. Na \olref[banach]{sec}, mencionamos que as
``peças'' da esfera que se obtêm não podem ser \emph{mensuráveis},
mas têm de ser dispersões infinitas impossíveis de visualizar. O mesmo vale para o nosso
uso da Escolha ao obter \olref{pseudobanachtarski}. E tudo isto merece explicação.

De novo, devemos esboçar alguma base (mas isto é \emph{apenas} um
esboço; talvez queira consultar uma entrada de manual sobre \emph{medida}).
Definir uma medida para um conjunto $X$ é atribuir um valor $\mu(E) \in
\Real$ a cada $E$ nalguma ``$\sigma$-álgebra'' sobre $X$. Os detalhes aqui
não são essenciais, exceto que a função $\mu$ deve obedecer ao
princípio da aditividade enumerável: a medida de uma união enumerável de
conjuntos disjuntos é a soma das medidas individuais, isto é,
$\mu(\bigcup_{n < \omega} X_n) = \sum_{n < \omega}\mu(X_n)$ sempre que
os $X_n$ são disjuntos. Dizer que um conjunto é ``não mensurável'' é
dizer que nenhuma medida pode ser atribuída de modo adequado. Agora, usando o nosso
$\rotationsgroup$ de antes:

\begin{cor}[Vitali]
Seja $\mu$ uma medida tal que $\mu(\onesphere) = 1$, e tal que
$\mu(X) = \mu(Y)$ se $X$ e $Y$ são congruentes. Então
$\funimage{\rho}{C}$ é não mensurável para todo o $\rho \in
\rotationsgroup$. 
\end{cor}

\begin{proof}
Por redução ao absurso, suponhamos o contrário. Seja então $\mu(\funimage{\sigma}{C}) =
r$ para algum $\sigma \in \rotationsgroup$ e algum $r \in \Real$. Para
qualquer $\rho \in \rotationsgroup$, $\funimage{\rho}{C}$ e $\funimage{\sigma}{C}$ são
congruentes, logo  $\mu(\funimage{\rho}{C}) = r$ para qualquer $\rho \in
\rotationsgroup$. Por \olref{vitalicover} e \olref{vitalinooverlap}, $\onesphere =
\bigcup_{\rho \in \rotationsgroup}\funimage{\rho}{C}$ é uma união enumerável de conjuntos disjuntos dois a dois. Logo a aditividade enumerável dita que
$\mu(\onesphere) = 1$ é a soma das medidas de cada
$\funimage{\rho}{C}$, isto é,
\[
	1 = \mu(\onesphere) = \sum_{\rho \in \rotationsgroup}\mu(\funimage{\rho}{C}) = \sum_{\rho \in \rotationsgroup}r
\]
Mas se $r = 0$ então $\sum_{\rho \in \rotationsgroup}r = 0$, e se $r
> 0$ então $\sum_{\rho \in \rotationsgroup}r = \infty$. 
\end{proof}

\end{document}
